\documentclass[aspectratio=169]{beamer}
\usetheme{SimplePlus}

\usepackage{tikz}
\usepackage{amsmath}
\usepackage{amssymb}
\usepackage{bm}
\usepackage{mathtools}
\DeclareMathOperator{\grad}{\nabla}


\title[Weak Formulations \& Lax-Milgram]{Weak Formulations and the Lax-Milgram Lemma}
\author{Nicol\'as A. Barnafi}
\date{}

\begin{document}

\begin{frame}
    \titlepage
\end{frame}

\begin{frame}{Outline}
    \tableofcontents
\end{frame}


\section{Weak Formulation of the Laplace Operator}

\begin{frame}{The Strong Form and Integration by Parts}
    For a bounded Lipschitz domain $\Omega$, \emph{in the sense of distributions}:
    $$ -\Delta u = f \quad \text{in } H^{-1}(\Omega) $$

    \vspace{0.2cm}
    \textbf{Weak formulation:} Test the equation against an arbitrary smooth test function $v$:
    $$ -\langle\Delta u, v \rangle = \langle f, v\rangle  $$

    \vspace{0.2cm}
    Distributional defition (i.e. integration by parts) yields:
    $$ -\langle \gamma_N \nabla u,\gamma_0 v\rangle + (\nabla u, \nabla v) = \langle f, v \rangle $$
    where $\gamma_0(u) = u|_{\partial\Omega}$ is the Dirichlet trace and $\gamma_N(\vec v)=(\vec v\cdot \vec n)|_{\partial\Omega}$ is the normal trace.
\end{frame}

\begin{frame}{Natural vs. Essential Boundary Conditions}
    The boundary integral reveals the natural decomposition of boundary conditions:
    $$ \int_{\partial\Omega} (\gamma_0 v) (\gamma_N \nabla u) \, ds = \int_{\Gamma_D} (\gamma_0 v) (\gamma_N \nabla u) \, ds + \int_{\Gamma_N} (\gamma_0 v) (\gamma_N \nabla u) \, ds $$

    \begin{itemize}
        \item \textbf{Essential (Dirichlet) Conditions:} Imposed directly on the solution space. 

            We define the test space such that $\gamma_0 v = 0$ on $\Gamma_D$. 
        \item \textbf{Natural (Neumann) Conditions:} Imposed via the formulation itself.

            If $\nabla u \cdot n = h$ on $\Gamma_N$, we substitute this directly into the boundary integral.
    \end{itemize}

$$ \langle \gamma_N \nabla u, \gamma_D v \rangle \quad\to\quad \langle h, \gamma_D v \rangle_{H^{-1/2}(\Gamma_N),H^{1/2}(\Gamma_N)} $$
\end{frame}

\begin{frame}{The Weak Formulation}
    Let $\Gamma_D = \partial\Omega$ (pure homogeneous Dirichlet problem for simplicity). We define our solution and test space as $V_0 = H^1_0(\Omega)$.

    \vspace{0.3cm}
    \textbf{The Weak Problem:}
    Find $u \in V_0$ such that
    $$ (\nabla u, \nabla v)_{0,\Omega} = \langle f, v \rangle \quad \forall v \in V_0 $$

    \vspace{0.3cm}
    \textbf{Observations:}
    \begin{itemize}
        \item The Laplacian is now interpreted distributionally.
        \item The solution and test spaces are identical ($W = V = H^1_0(\Omega)$).
        \item The Dirichlet boundary condition does not appear in the integral equation.
    \end{itemize}
\end{frame}

\section{Abstract Problems and Well-Posedness}

\begin{frame}{The Abstract Variational Problem}
    Let $(W,\|\cdot\|_W)$, $(V, \|\cdot \|_V)$ normed vector spaces. Generally, $V$ is reflexive Banach space and $W$ is Banach.

    \vspace{0.3cm}
    \textbf{The Continuous Problem:}
    Find $u \in W$ such that
    $$ a(u, v) = f(v) \quad \forall v \in V $$

    \vspace{0.3cm}
    \textbf{Components:}
    \begin{itemize}
        \item $a(\cdot, \cdot)$: Continuous bilinear form on $W \times V$, i.e. $a \in \mathcal{L}(W \times V; \mathbb{R})$.
        \item $f(\cdot)$: Continuous linear form on $V$, i.e. $f \in V'$.
    \end{itemize}
\end{frame}

\begin{frame}{Hadamard Well-Posedness}
    To consider solving a boundary value problem, we must \emph{good} problems.

    \vspace{0.3cm}
    \textbf{Definition (Well-posedness in the sense of Hadamard):}
    The abstract problem is well-posed if it satisfies two conditions:
    \begin{enumerate}
        \item It admits a unique solution $u \in W$.
        \item It satisfies the following \textit{a priori} or \emph{stability} estimate:
        $$ \exists c > 0, \forall f \in V', \quad \|u\|_W \le c\|f\|_{V'} $$
    \end{enumerate}

    \vspace{0.3cm}
    \idea{ The solution exists, is unique, and depends continuously on the given data $f$.}
\end{frame}


\section{The Lax-Milgram Lemma}

\begin{frame}{Abstract Framework and Ellipticity}
    We can frame the Poisson weak formulation as an abstract problem $a(u,v) = f(v)$ by defining the bilinear form $a(u,v) := (\nabla u, \nabla v)_{0,\Omega}$.

    \vspace{0.2cm}
    \textbf{Definition (Boundedness):}
    A bilinear form $a(\cdot, \cdot)$ on $H \times H$ is bounded (continuous) if:

    There exists $C > 0$ such that:
    $$ |a(u, v)| \le C \|u\|_H \|v\|_H \quad \forall u, v \in H $$

    \vspace{0.2cm}
    \textbf{Definition (Ellipticity):}
    A bilinear form $a(\cdot, \cdot)$ defined on a Hilbert space $H$ is elliptic if: 

    There exists a constant $\alpha > 0$ such that:
    $$ a(x, x) \ge \alpha \|x\|^2_H \quad \forall x \in H $$

\end{frame}

\begin{frame}{The Lax-Milgram Lemma}
    \begin{block}{Lemma (Lax-Milgram)}
    Consider: 
    \begin{itemize}
        \item $H$ a Hilbert space. 
        \item A bounded, elliptic bilinear form $a : H \times H \to \mathbb{R}$ with constants $C$ and $\alpha$. 
        \item A linear functional $f \in H'$.
    \end{itemize}

        Then, the abstract problem is well-posed: There exists a unique solution $u \in H$ such that
        $$ a(u, v) = f(v) \quad \forall v \in H, $$

        And the inverse is stable:
        $$ \|u\|_H \le \frac{1}{\alpha}\|f\|_{H'} $$
\end{block}
\end{frame}


\begin{frame}{Application to Poisson problem}
    \begin{itemize}
        \item $H\coloneqq (H_0^1(\Omega), (\nabla \cdot, \nabla \cdot))$ is a Hilbert space.
        \item The linear operator $f$ is bounded: 
            $$ f(v) = \langle f, v\rangle \leq \| f \|_{H^{-1}(\Omega)} \| v \|_H $$
        \item The bilinear form $a$ is bounded:
            $$ a(u,v) = (\nabla u, \nabla v) \leq \| \nabla u \| \| \nabla v \| $$
        \item The bilinear form $a$ is elliptic:
            $$ a(v,v) = (\nabla v, \nabla v) = \| v \|_H $$
    \end{itemize}
    Note that $H$ is Hilbert because of Poincaré, and we can repeat proof using $(H_0^1(\Omega), (\cdot, \cdot) + (\nabla \cdot, \nabla \cdot)) $. 
\end{frame}

\section{Proof of the Lax-Milgram Lemma}

\begin{frame}{Step 1: Induced Operators and the Riesz Map}
    Every continuous bilinear form $a: H \times H \to \mathbb{R}$ induces a bounded linear operator $A: H \to H'$ defined by:
    $$ \langle Au, v \rangle_{H', H} = a(u, v) \quad \forall v \in H $$

    By the Riesz Representation Theorem, there exists the Riesz map $R: H \to H'$.

    \vspace{0.3cm}
    We can pull the operator back into $H$ by defining $R^{-1} \circ A : H \to H$. The abstract problem $a(u,v) = f(v)$ is thus equivalent to the operator equation:
    $$ Au = f \text{ in } H' \iff (R^{-1} \circ A)u \coloneqq \mathbb A u = R^{-1}f \coloneqq F \text{ in } H $$
\end{frame}

\begin{frame}{Step 2: Fixed Point Formulation}
        $$ \mathbb A u = F $$
    To prove existence and uniqueness, we construct a fixed-point mapping.

    For a parameter $\rho > 0$, define the map $T: H \to H$ as:
    $$ T(u) = u - \rho (\mathbb Au - F) $$

    Clearly, $u$ is a fixed point of $T$ iff $Au = f$. 

    \idea{We impose contraction through $\rho$}
\end{frame}

\begin{frame}{Step 3: Contraction Mapping Argument}
    For $u, v \in H$, analyze the distance between their images:
    \begin{align*}
        \|T(u) - T(v)\|_H^2 &= \|u - v - \rho R^{-1}A(u - v)\|_H^2 \\
        &= \|u - v\|_H^2 - 2\rho(u - v, R^{-1}A(u - v))_H + \rho^2\|R^{-1}A(u - v)\|_H^2
    \end{align*}

    Using the Riesz map properties, $(x, R^{-1}y)_H = \langle y, x \rangle_{H',H}$:
    $$ \|T(u) - T(v)\|_H^2 = \|u - v\|_H^2 - 2\rho\langle A(u - v), u - v \rangle_{H',H} + \rho^2\|A(u - v)\|_{H'}^2 $$

    Applying ellipticity ($\langle Au, u \rangle \ge \alpha\|u\|^2$) and boundedness ($\|Au\|_{H'} \le C\|u\|$):
    $$ \|T(u) - T(v)\|_H^2 \le (1 - 2\rho\alpha + \rho^2 C^2)\|u - v\|_H^2 $$

    $T$ is a contraction if $1 - 2\rho\alpha + \rho^2 C^2 < 1$, which holds for $\rho \in (0, 2\alpha/C^2)$. Banach's Fixed Point Theorem guarantees a unique solution.
\end{frame}

\begin{frame}{Step 4: Stability Estimate}
    Finally, we derive the \textit{a priori}.

    By ellipticity, we have:
    $$ \alpha\|u\|_H^2 \le a(u, u) $$

    Since $u$ is the solution, $a(u,u) = f(u)$:
    $$ f(u) = \langle f, u \rangle \le \|f\|_{H'}\|u\|_H $$

    Then:
    $$ \alpha\|u\|_H^2 \le \|f\|_{H'}\|u\|_H $$

    Dividing by $\alpha\|u\|_H$ (assuming $u \neq 0$) gives the required stability bound:
    $$ \|u\|_H \le \frac{1}{\alpha}\|f\|_{H'} $$
    \qed
\end{frame}

\begin{frame}
    \maketitle
\end{frame}
\end{document}
